% =====================================================================
%  Formulas for the Gesture MIDI Controller paper (methodology section)
%  Each block is self-contained; paste the ones you need.
%  Requires: \usepackage{amsmath}
% =====================================================================

% ---------------------------------------------------------------------
% 1. HSV colour segmentation (with red hue-wrap)
% ---------------------------------------------------------------------
% A pixel p is part of marker m's mask if its HSV value lies within the
% calibrated lower/upper bounds. OpenCV hue is in [0,179]; red wraps at 0,
% so its range is expressed as the union of two intervals.
\begin{equation}
M_m(p) =
\begin{cases}
\mathbf{1}\!\left[\, c^{\text{lo}}_m \preceq \mathrm{HSV}(p) \preceq c^{\text{hi}}_m \,\right],
  & h^{\text{lo}}_m \le h^{\text{hi}}_m \\[4pt]
\mathbf{1}\!\left[ \mathrm{HSV}(p) \in R_1 \cup R_2 \right],
  & h^{\text{lo}}_m > h^{\text{hi}}_m
\end{cases}
\label{eq:segmentation}
\end{equation}
% where for the wrap case R1 = [0,h^hi] and R2 = [h^lo,179] in hue,
% sharing the same S,V bounds.

% ---------------------------------------------------------------------
% 2. Calibration: per-channel HSV bounds from samples
% ---------------------------------------------------------------------
% Bounds are mean +/- k.sigma over N sampled pixels, with capped floors
% on saturation and value for lighting tolerance.
\begin{align}
\mu_c &= \tfrac{1}{N}\sum_{i=1}^{N} x_{i,c}, \qquad
\sigma_c = \sqrt{\tfrac{1}{N}\sum_{i=1}^{N}(x_{i,c}-\mu_c)^2} \\
h^{\text{lo}} &= \mu_H - 3\sigma_H, \quad h^{\text{hi}} = \mu_H + 3\sigma_H \\
s^{\text{lo}} &= \mathrm{clip}(\mu_S - 3\sigma_S,\; 30,\; 150) \\
v^{\text{lo}} &= \mathrm{clip}(\mu_V - 5\sigma_V,\; 20,\; 110)
\label{eq:calibration}
\end{align}

% ---------------------------------------------------------------------
% 3. Marker centroid via image moments
% ---------------------------------------------------------------------
% Spatial moments of the (binary) mask give the marker centroid.
% Fingertips use the largest contour; the wrist loop uses the whole mask.
\begin{equation}
M_{ji} = \sum_{x}\sum_{y} x^{j}\, y^{i}\, B(x,y), \qquad
(\bar{x},\bar{y}) = \left( \frac{M_{10}}{M_{00}},\; \frac{M_{01}}{M_{00}} \right)
\label{eq:centroid}
\end{equation}

% ---------------------------------------------------------------------
% 4. Exponential moving-average (EMA) smoothing
% ---------------------------------------------------------------------
% Applied to coordinates, pinch and tilt (alpha=0.3) and to the wrist
% centroid (alpha=0.4). The step response reaches (1-(1-alpha)^{n+1}).
\begin{equation}
s_t = \alpha\, x_t + (1-\alpha)\, s_{t-1}, \qquad \alpha \in (0,1]
\label{eq:ema}
\end{equation}

% ---------------------------------------------------------------------
% 5. Pinch gesture (normalized thumb-index distance)
% ---------------------------------------------------------------------
\begin{equation}
d = \left\lVert (i_x,i_y) - (t_x,t_y) \right\rVert_2, \qquad
p = \min\!\left(1,\; \max\!\left(0,\; \frac{d}{0.35}\right)\right)
\label{eq:pinch}
\end{equation}

% ---------------------------------------------------------------------
% 6. Tilt (wrist-to-index orientation)
% ---------------------------------------------------------------------
\begin{equation}
\theta = \frac{1}{\pi/2}\,
\operatorname{atan2}\!\big(i_x - w_x,\; (1-i_y) - (1-w_y)\big),
\qquad \theta \in [-1,1]
\label{eq:tilt}
\end{equation}

% ---------------------------------------------------------------------
% 7. Gesture -> MIDI mapping (absolute and relative)
% ---------------------------------------------------------------------
% Absolute: a normalized gesture g is quantized to a 7-bit CC.
\begin{equation}
\mathrm{CC}^{\text{abs}} = \left\lfloor \mathrm{clip}(g,0,1)\cdot 127 \right\rfloor
\label{eq:abs}
\end{equation}
% Relative: a signed per-frame delta is accumulated and clamped (v0 = 64).
\begin{equation}
v_t = \mathrm{clip}\!\big(v_{t-1} + \lfloor \delta_t \cdot 254 \rfloor,\; 0,\; 127\big)
\label{eq:rel}
\end{equation}

% ---------------------------------------------------------------------
% 8. End-to-end per-frame latency model
% ---------------------------------------------------------------------
\begin{equation}
T_{\text{frame}} = T_{\text{capture}} + T_{\text{hsv}}
                 + T_{\text{search}} + T_{\text{map}},
\qquad
\mathrm{FPS}_{\max} = \frac{1}{T_{\text{frame}}}
\label{eq:latency}
\end{equation}

% ---------------------------------------------------------------------
% 9. MediaPipe duty cycle (worker runs every N-th frame)
% ---------------------------------------------------------------------
% With N=5, MediaPipe touches only 1/N of frames and runs off the hot path,
% so the amortized per-frame tracking cost is bounded by the HSV path.
\begin{equation}
\bar{T}_{\text{track}} \approx T_{\text{hsv+search}}
   + \tfrac{1}{N}\, T_{\text{MP}}\big|_{\text{(off hot path)}}
\label{eq:duty}
\end{equation}
